\section{Endpoint noise tails and the square-root scale}
\label{sec:square-root-tail}
This section removes the logarithmic loss in the preceding sufficient
condition and determines the exact vanishing scale of the operator norm
used in its proof. The two conclusions have different scopes: sharpness
for an operator norm is not necessity for nonnegative weighting densities.
All norms below use uniform probability measure on the Boolean cube.
Fix $0<\rho<1$, put $q_0=1+\rho$, $p_0=q_0/(q_0-1)=1+1/\rho$,
and write $Q=P_{\sqrt\rho}$. Let $H_m$ be the orthogonal Walsh projection
onto degrees at least $m$, with $H_m=0$ if $m>L$.

\begin{theorem}[Endpoint tail bound and exact vanishing scale]
\label{thm:square-root-tail}
There is $C_\rho>0$, depending only on the fixed parameter $\rho$, such that,
for $1\leq m\leq L$,
\begin{equation}\label{eq:327-tail-upper}
 N_L(m):=\|QH_m\|_{2\to p_0}
 \leq \rho^{m/4}+e^{-cm}+\exp(-C_\rho c^2m^2/L),
 \qquad c=\frac{(-\log\rho)(q_0-1)}{8(1-q_0/2)}.
\end{equation}
For every sequence of positive integers $m_L$,
\[
 N_L(m_L)\longrightarrow0
 \quad\Longleftrightarrow\quad m_L/\sqrt L\longrightarrow\infty.
\]
More precisely, if $m_L/\sqrt L\to C\in[0,\infty)$, then
\begin{equation}\label{eq:327-tail-lower}
 \liminf_{L\to\infty}N_L(m_L)
 \geq \exp\!\left(-\frac{1-\rho^2}{12\rho}C^2\right).
\end{equation}
In particular $N_L(m_L)\to1$ when $m_L=o(\sqrt L)$.
No equality for the positive-$C$ transition profile is asserted.
\end{theorem}

We first establish the quantitative modulus needed for the upper bound.
Our external input is the distribution-sensitive hypercontractive theorem
of Kirshner and Samorodnitsky, Theorem~1.8 and Lemma~2.14 in
\url{https://arxiv.org/html/1909.11929v1}. With
$\epsilon=(1-\sqrt\rho)/2$ and
$x=L^{-1}\log_2(\|f\|_{q_0}/\|f\|_1)$, their theorem gives
\begin{equation}\label{eq:327-KS}
 \|Qf\|_2\leq 2^{L\eta(x,\epsilon)}\|f\|_{q_0}.
\end{equation}
Its argument belongs to $[0,(q_0-1)/q_0]$. Their function is continuous
there and strictly negative for positive argument. We use this attributed
inequality; we do not claim a new proof of that theorem.

\begin{lemma}[Quadratic modulus of the concentration deficit]
\label{lem:327-modulus}
For the fixed parameters above there is $C_\rho>0$ such that
\[
 -\log(2)\eta(x,\epsilon)\geq C_\rho(\log(2)x)^2
 \qquad (0\leq x\leq(q_0-1)/q_0).
\]
The continuous quotient at zero has value
$(2-q_0)/(3q_0(q_0-1))>0$. The global constant is not claimed to equal
this local value, or to be uniform as $\rho$ approaches either endpoint.
\end{lemma}
\begin{proof}
Write $d=(1-\rho)/2=2\epsilon(1-\epsilon)$ and consider signs $(X,Y)$
with probabilities $q_0/4$ on each equal pair and $(2-q_0)/4$ on each
unequal pair. All four probabilities are positive. For
$\sigma=(1-t)/2$, any joint table with both negative-sign marginals
equal to $\sigma$ has entries
\[
 \nu(--)=\sigma-z,\quad \nu(-+)=\nu(+-)=z,\quad
 \nu(++)=1-\sigma-z,\qquad 0\leq z\leq\sigma.
\]
For $t\geq0$, its base-two entropy is
$H_2(\sigma)+\sigma H_2(z/\sigma)+(1-\sigma)H_2(z/(1-\sigma))$,
with boundary values understood by continuity. Thus
\begin{align*}
 -D_2(\nu\Vert P)={}&H_2(\sigma)-1
   +\sigma H_2(z/\sigma)+(1-\sigma)H_2(z/(1-\sigma))\\
 &+(1-2z)\log_2(1-d)+2z\log_2d.
\end{align*}
Maximizing this expression over $z$ is exactly the definition of
$\phi(\sigma,d)$ in Section~2.1.6 of the cited source. Consequently
$\phi(\sigma,d)=-I(t,t)/\log2$, where $I(t,t)$ is the minimum natural
relative entropy subject to the two mean constraints. This is an identity
between finite probability tables, requiring no asymptotic identification
of a large-deviation rate.

Near zero, the covariance matrix
$\left(\begin{smallmatrix}1&\rho\\\rho&1\end{smallmatrix}\right)$
is positive definite. The analytic inverse function theorem applied to
the gradient of the joint log moment-generating function produces a unique
small tilt realizing means $(t,t)$. If $P_\lambda$ denotes that tilt,
the identity
\[
 D(\nu\Vert P)=D(\nu\Vert P_\lambda)
  +\lambda_1t+\lambda_2t-\log\mathbb E_Pe^{\lambda_1X+\lambda_2Y}
\]
proves that it minimizes entropy. Swap symmetry and local uniqueness
give $\lambda_1=\lambda_2=\lambda$. Therefore, with natural logarithms,
\[
 K(\lambda)=\log\left(\frac{q_0}{2}\cosh(2\lambda)
                          +\frac{2-q_0}{2}\right),\qquad
 I(t,t)=2t\lambda-K(\lambda),\quad K'(\lambda)=2t.
\]
Writing $b=q_0/3-q_0^2/2$, analytic expansion and inversion give
\begin{align*}
 K(\lambda)&=q_0\lambda^2+b\lambda^4+O(\lambda^6),\\
 \lambda&=t/q_0-2bt^3/q_0^4+O(t^5),\\
 I(t,t)&=t^2/q_0+(q_0/2-1/3)t^4/q_0^3+O(t^6).
\end{align*}
Also $I_1(t)=\log2[1-H_2((1-t)/2)]=t^2/2+t^4/12+O(t^6)$.
The definition of $\eta$ in Section~2.1.7 of the cited source has
$1-H_2(\sigma)=q_0x/(q_0-1)$. Hence, setting $u=(\log2)x$,
\[
 u=\frac{q_0-1}{q_0}I_1(t),\qquad
 -\log2\,\eta(x,\epsilon)=\tfrac12 I(t,t)-I_1(t)/q_0.
\]
The first relation has an analytic inverse in $t^2$ near zero. Substitution
in the second gives
\[
 -\log2\,\eta(x,\epsilon)
 =\frac{(q_0-1)(2-q_0)}{12q_0^3}t^4+O(t^6)
 =\frac{2-q_0}{3q_0(q_0-1)}u^2+O(u^3).
\]
Continuity and strict negativity of $\eta$ away from zero, supplied by
the cited definition and lemma, now make the quotient a positive
continuous function on a compact interval. Its positive minimum proves
the assertion.
\end{proof}

\begin{proof}[Proof of Theorem~\ref{thm:square-root-tail}]
Self-adjointness, commutation and finite-dimensional duality give
$N_L(m)=\|H_mQ\|_{q_0\to2}$. Normalize $\|f\|_{q_0}=1$ and put
$a=(-\log\rho)/[4(1-q_0/2)]$, $M=e^{am}$, so $c=a(q_0-1)/2$.
Decompose $f=b+t$ according to $|f|\leq M$ and $|f|>M$.
The bounded part satisfies $\|b\|_2\leq M^{1-q_0/2}$, whence
\[
 \|H_mQb\|_2\leq\rho^{m/2}M^{1-q_0/2}=\rho^{m/4}.
\]
Here the projection is used only in $L^2$, with its spectral multiplier;
no $L^{p_0}$ contraction of a Fourier cutoff is assumed.

For the tail, $\|t\|_1\leq e^{-a(q_0-1)m}$ and $\|t\|_{q_0}\leq1$.
If $\|t\|_{q_0}\leq e^{-cm}$, ordinary endpoint hypercontractivity
bounds $\|H_mQt\|_2$ by $e^{-cm}$. Otherwise
$\log(\|t\|_{q_0}/\|t\|_1)\geq cm$, and
\eqref{eq:327-KS} with Lemma~\ref{lem:327-modulus} bounds it by
$\exp(-C_\rho c^2m^2/L)$.
For signed functions, the positive kernel gives
$|Qt|\leq Q|t|$, so the same refined estimate applies. A zero tail is
harmless; a case forcing a concentration argument above its admissible
interval is empty. Adding the two possible tail bounds yields
\eqref{eq:327-tail-upper}. If $m_L/\sqrt L\to\infty$, each displayed
term vanishes on indices $m_L\leq L$, while the other indices give
the zero operator.

For the converse, for $0<t<1$ define the unit $L^2$ product
\[
 f_{L,t}(x)=\prod_{i=1}^L\frac{1+tx_i}{\sqrt{1+t^2}}.
\]
Its Walsh coefficient at $A$ is $t^{|A|}/(1+t^2)^{L/2}$, so exactly
\[
 \|(I-H_m)f_{L,t}\|_2^2
 =\Pr\{\operatorname{Bin}(L,t^2/(1+t^2))<m\}.
\]
Ordinary endpoint hypercontractivity and the triangle inequality imply
\begin{equation}\label{eq:327-witness}
 N_L(m)\geq\|Qf_{L,t}\|_{p_0}
 -\sqrt{\Pr\{\operatorname{Bin}(L,t^2/(1+t^2))<m\}}.
\end{equation}
In particular the projection, not the original positive product, is the
high-degree witness. It may be signed.

Product factorization gives $\log\|Qf_{L,t}\|_{p_0}=LF_\rho(t)$, where
\[
 F_\rho(t)=\frac1{p_0}\log\frac{(1+\sqrt\rho\,t)^{p_0}
                                  +(1-\sqrt\rho\,t)^{p_0}}2
                  -\frac12\log(1+t^2).
\]
Both bases are positive near zero, so this is even analytic also for
noninteger $p_0$. Expanding
$\mathbb E(1+zX)^{p_0}=1+A z^2+Bz^4+O(z^6)$ with
$A=p_0(p_0-1)/2$ and $B=p_0(p_0-1)(p_0-2)(p_0-3)/24$ gives
\[
 \frac1{p_0}\log\mathbb E(1+zX)^{p_0}
 =\frac{p_0-1}{2}z^2
 -\frac{(p_0-1)(p_0^2+p_0-3)}{12}z^4+O(z^6).
\]
The quadratic terms cancel since $(p_0-1)\rho=1$, and thus
$F_\rho(t)=-d_\rho t^4+O_\rho(t^6)$, where
$d_\rho=(1-\rho^2)/(12\rho)>0$.

Suppose $m_L/\sqrt L\to C<\infty$. Fix $A>C$ and set
$t_L^2=A/\sqrt L$. Then $Lt_L^4=A^2$ and $Lt_L^6\to0$, so the
first term in \eqref{eq:327-witness} tends to $e^{-d_\rho A^2}$.
The binomial mean is $\mu_L=A\sqrt L/(1+A/\sqrt L)$ and its
variance is at most $\mu_L$. For large $L$,
$\mu_L-m_L\geq(A-C)\sqrt L/2$; Chebyshev's inequality bounds
the low-tail probability by $\mu_L/(\mu_L-m_L)^2=O_{A,C}(L^{-1/2})$.
Taking the lower limit in \eqref{eq:327-witness}, then the supremum over
fixed $A>C$, proves \eqref{eq:327-tail-lower}. There is no interchange
of nonuniform limits. At $C=0$, the endpoint upper bound $N_L(m)\leq1$
proves convergence to one.
Finally, if $m_L/\sqrt L$ fails to tend to infinity, it has a bounded
subsequence and then a subsequence converging to a finite $C$.
The positive lower bound on that subsequence precludes $N_L(m_L)\to0$.
\end{proof}

\begin{corollary}[Marginal rigidity without the logarithmic loss]
\label{cor:327-marginals}
Let $r_L\geq0$, $\mathbb E r_L=1$, and
$\|r_L\|_s\leq B$ for $s=(1+\rho)/(1-\rho)$ and fixed finite $B$.
Assume all weighting-law marginals on at most $k_L$ coordinates are
uniform. If $k_L/\sqrt L\to\infty$, then
\[
 \|S_{r_L}-P_\rho\|_{2\to2}\to0,\qquad
 \lambda_1(T_{r_L})\to1,\quad\lambda_2(T_{r_L})\to\rho,
 \qquad \gap(T_{r_L})\to1-\rho.
\]
The norm conclusion concerns $S_{r_L}=QM_{r_L}Q$; it does not assert
$\|T_{r_L}-P_\rho\|\to0$ for $T_{r_L}=M_{\sqrt{r_L}}P_\rho M_{\sqrt{r_L}}$.
\end{corollary}
\begin{proof}
Put $d=\lfloor k_L/2\rfloor$, let $\Pi$ project onto degrees at most
$d$, and let $H=I-\Pi=H_{d+1}$. With $h=r_L-1$ and $D=QM_hQ$,
exact matching gives $\Pi D\Pi=0$: the Walsh entry uses
$\widehat h(A\mathbin\triangle C)$ with $|A\mathbin\triangle C|\leq2d\leq k_L$.
Consequently $D=HD+\Pi DH$. Since $1/s+2/p_0=1$,
H\"older's inequality and endpoint hypercontractivity bound each block
by $(B+1)N_L(d+1)$; for the projected low factor use
$\|Q\Pi f\|_{p_0}\leq\|\Pi f\|_2\leq\|f\|_2$.
Thus $\|D\|\leq2(B+1)N_L(d+1)\to0$. If $k_L=L$, exact matching
already forces $r_L=1$. Self-adjoint eigenvalue perturbation yields the
top two limits for $S_{r_L}$. With $A=QM_{\sqrt{r_L}}$ we have
$S_{r_L}=AA^*$ and $T_{r_L}=A^*A$; their full finite spectra agree,
including zeros. This proves the asserted eigenvalue and normalized-gap
limits even when the densities vanish.
\end{proof}

\paragraph{Scope and relation to earlier results.}
The new sufficient condition supersedes the logarithmic scale of
Theorem~\ref{thm:quadratic-marginal-rigidity}; that earlier bound remains
valid. Together with Theorem~\ref{thm:growing-collapse}, it leaves a
genuine density question near the square-root scale and in the gap
between fixed sub-square-root powers and that scale. The signed witness
in \eqref{eq:327-witness} does not supply a nonnegative mean-one density
with exact marginals and a collapsing transfer gap. Nor does equality
of spectra of $S_r,T_r$ give equality of their distances to $P_\rho$.
For example, $r=1+t\chi_{[L]}$, $0<t<1$, has uniform proper marginals
and bounded critical norm. Writing $\sqrt r=a+b\chi_{[L]}$ gives
$2ab=t$, so the $\{1,\chi_{[L]}\}$ off-diagonal entry of
$T_r-P_\rho$ is $t(1+\rho^L)/2\geq t/2$, whereas
$\|S_r-P_\rho\|=t\rho^{L/2}\to0$.
These distinctions correct the ambiguous norm wording of the previous
version without changing its formally stated $S_r$ theorem.

Endpoint hypercontractivity, Fourier projections, entropy tilting and
product norm factorization are classical inputs. The refined inequality
is attributed explicitly above. Existing tail-space literature also
contains square-root phenomena; this fact alone neither identifies nor
excludes an antecedent of this precise endpoint operator theorem.
Historical originality remains unestablished. These are finite-cube
results with global weights, not fixed-range Gibbs, local-dynamics,
physical-gap or continuum results.
